\section{Proposed Method}
\label{sec:method}

We formulate VQA as conditional answer generation: given an image $I$ and a Vietnamese
question $q$, the model generates a short free-form answer $y=(y_1,\dots,y_T)$.
ViBridge-VQA builds on Vintern-1B~\cite{doan2024vintern1b} and keeps all of its
pre-trained weights frozen. Instead of adapting the model, it changes what the decoder sees. 

\begin{figure}[t]
\centering
\resizebox{\linewidth}{!}{%
\begin{tikzpicture}[font=\small, node distance=3.5mm,
  frozen/.style={draw=gray!70, rounded corners=2pt, fill=gray!10, minimum height=9mm, align=center},
  train/.style={draw=green!45!black, rounded corners=2pt, fill=green!12, thick, minimum height=9mm, align=center},
  data/.style={align=center, text=black!75},
  arr/.style={-{Latex[length=1.6mm]}, thick, black!70},
  lbl/.style={font=\footnotesize\bfseries, anchor=west}]
% ---- row 1: official fine-tuning of Vintern-1B
\node[lbl] (t1) at (0,1.2) {(a) Vintern-1B, official fine-tuning};
\node[data] (img1) at (0.7,0) {image\\$\le$6 tiles + thumbnail\\($448^2$ each)};
\node[frozen, right=of img1] (vit1) {InternViT-300M};
\node[frozen, right=of vit1] (mlp1) {projector\\(patches, $2{\times}2$ merge)};
\node[data, right=of mlp1] (tok1) {256 tokens / tile\\($\le$1{,}792)};
\node[frozen, right=of tok1, text width=27mm] (llm1) {Qwen2.5-0.5B\\+ \textcolor{green!45!black}{\textbf{LoRA} $r{=}16$}};
\node[data, right=of llm1] (ans1) {answer};
\draw[arr] (img1) -- (vit1); \draw[arr] (vit1) -- (mlp1); \draw[arr] (mlp1) -- (tok1);
\draw[arr] (tok1) -- (llm1); \draw[arr] (llm1) -- (ans1);
% ---- row 2: ViBridge-VQA
\node[lbl] (t2) at (0,-1.5) {(b) ViBridge-VQA (ours)};
\node[data] (img2) at (0.7,-3.1) {image\\1 view\\($336^2$)};
\node[frozen] (vit2) at (vit1 |- img2) {InternViT-300M};
\node[train] (br) at ($(mlp1 |- img2)+(0,0.75)$) {\textbf{global bridge}\\(\texttt{[CLS]} $\to g$ tokens)};
\node[frozen] (pj) at ($(mlp1 |- img2)+(0,-0.75)$) {projector of Vintern-1B\\+ avg.\ pool to $k$};
\node[data] (tg) at (tok1 |- br) {$g$ global tokens\\(prefix)};
\node[data] (tl) at (tok1 |- pj) {$k$ local tokens\\(image slot)};
\node[frozen, text width=27mm] (llm2) at (llm1 |- img2) {Qwen2.5-0.5B\\(frozen, no LoRA)};
\node[data, right=of llm2] (ans2) {answer};
\draw[arr] (img2) -- (vit2);
\draw[arr] (vit2.east) -- ++(0.15,0) |- (br.west);
\draw[arr] (vit2.east) -- ++(0.15,0) |- (pj.west);
\draw[arr] (br) -- (tg); \draw[arr] (pj) -- (tl);
\draw[arr] (tg.east) -- ++(0.15,0) |- (llm2.west);
\draw[arr] (tl.east) -- ++(0.15,0) |- (llm2.west);
\draw[arr] (llm2) -- (ans2);
% legend
\node[frozen, minimum height=4mm, minimum width=5mm] (lf) at (13.2,-4.6) {};
\node[anchor=west] at (lf.east) {frozen};
\node[train, minimum height=4mm, minimum width=5mm] (lt) at (15.0,-4.6) {};
\node[anchor=west] at (lt.east) {trained};
\end{tikzpicture}}
\caption{Overview of ViBridge-VQA compared with the official fine-tuning of Vintern-1B.
(a) The official recipe encodes up to seven views per image, passes up to 1{,}792 visual
tokens to the decoder and adapts the decoder with LoRA. (b) ViBridge-VQA encodes a single
view and passes $g$ global tokens, produced from the \texttt{[CLS]} vector by a trainable
bridge, and $k$ local tokens, produced by the frozen projector of Vintern-1B and
average-pooled. Only the global bridge is trained (1.35\% of the parameters for
$g=14$); the selected budget is $g=14$, $k=36$.}
\label{fig:arch}
\end{figure}

\subsection{Overall Architecture}
\label{sec:encoding}

As illustrated in Fig.~\ref{fig:arch}, ViBridge-VQA keeps the three components of
Vintern-1B: the InternViT-300M vision encoder (24 layers, width 1024), its two-layer MLP
projector $P$, and the Qwen2.5-0.5B-Instruct language decoder (24 layers, width 896).
Each image is resized to one $336\times336$ view ($24\times24$ patches of $14\times14$ pixels);
the encoder returns one vector per patch and a global \texttt{[CLS]} vector $v\in\mathbb{R}^{1024}$.
Vintern-1B uses only the patch vectors: it merges every $2\times2$ block of neighbouring
patches into one 4096-dimensional vector (pixel shuffle) and maps each merged vector into
the decoder input space with $P$, which was pre-trained to align the two backbones. At our
resolution this yields a $12\times12$ grid of 144 region tokens. ViBridge-VQA uses both outputs of the encoder.

\subsection{Global Tokens}
\label{sec:bridge}

The \texttt{[CLS]} vector summarises the whole scene, but Vintern-1B never passes it to
the decoder. A small trainable bridge $g_\phi$ expands it into $g$ tokens with two linear
heads:
\begin{equation}
  g_\phi(v) = \big[\,W_a v + b_a;\ \mathrm{reshape}_{(g-1)\times 896}(W_e v + b_e)\,\big]
  \in \mathbb{R}^{g\times 896}.
  \label{eq:bridge}
\end{equation}
where $W_a\in\mathbb{R}^{896\times1024}$ and $W_e\in\mathbb{R}^{(g-1)896\times1024}$. The first
head produces an \emph{anchor} token, and the second produces $g-1$ complementary tokens that
present the global content to the decoder from different learned directions. The bridge has
$1025\cdot896\cdot g$ parameters, i.e., 12{,}857{,}600 for $g=14$, and its $g$ tokens are
placed before the prompt. We use $g=14$ because, in a preliminary study, the F1 of this bridge
alone rose with $g$ up to 14 and then plateaued.

\subsection{Local Tokens}
\label{sec:local}

Global tokens alone cannot describe small objects, printed text or the position of
objects. ViBridge-VQA therefore also passes local tokens, taken directly from the 144
region tokens of the frozen projector $P$. To control their number, we average-pool the
$12\times12$ grid to an $s\times s$ grid, so that $k=s^2$ with $s\in\{1,3,6,12\}$: $k=144$
keeps every region token, and $k=36$ averages each $2\times2$ block. The pooled tokens
replace the embeddings of $k$ image-context placeholders inside Vintern-1B's own image
slot, \texttt{<img>}\,$\cdots$\,\texttt{</img>}, in the prompt. Because they come from the
pre-trained projector rather than from a newly learned projection of the patches, the
local tokens stay in the representation that the decoder was trained to read, and no
parameters are needed to produce them.

\subsection{Training Objective}
\label{sec:training}

Only the bridge parameters $\phi$ are trained; the vision encoder, the projector and the
decoder remain frozen, and no adapter is added to the decoder. We minimise the
cross-entropy of the answer tokens:
\begin{equation}
  \mathcal{L} = -\sum_{t=1}^{T} \log p\big(y_t \mid y_{<t},\, g_\phi(v),\, \mathrm{pool}_s(P(x)),\, q\big),
  \label{eq:ce}
\end{equation}
where $x$ denotes the merged patch vectors. ViBridge-VQA can also be combined with decoder adaptation: in that variant, a rank-16
LoRA~\cite{hu2022lora} on the query, key, value and output projections of every decoder layer is
trained jointly with the bridge, which adds 2{,}162{,}688 parameters. Throughout the paper,
percentages of trained parameters are relative to all parameters of the model, i.e., the
frozen components of Vintern-1B (about 0.94B, including the vision encoder, projector,
embeddings and output layer) plus the trained modules. With $g=14$, the bridge alone trains
12.86M parameters (1.35\%), and the bridge with LoRA 15.02M (1.57\%).

\subsection{Choosing the Token Budget}
\label{sec:budget}

The cost of the decoder grows with the number of visual tokens $g+k$, so we treat the pair
$(g,k)$ as a budget to be chosen. In our experiments, accuracy increases monotonically
with the number of visual tokens, so every configuration lies on the accuracy--cost Pareto front and
the front alone does not single one out. We therefore select its
\emph{knee}~\cite{satopaa2011kneedle}, the point after which additional tokens bring
diminishing returns. After normalising cost $c$ and accuracy $a$ to $[0,1]$ over the front,
the knee is the configuration that maximises $\hat a - \hat c$, i.e., the one lying
furthest above the line joining the cheapest and the most accurate configuration. The knee is computed on the validation split only, and we check the choice with the
multi-criteria method TOPSIS~\cite{hwang1981topsis}, also on validation, under several
weightings of accuracy and cost. Test results are computed for every configuration only to be
reported; no choice is based on them.
